\makeabstract
    {
    Visualizing Histone-Protamine Replacement in Spermiogenesis using Atomic Force Microscopy
    }
    {
    Renna Chang\textsuperscript{1}, Ashley Carter\textsuperscript{1}
    }
    {
    \textsuperscript{1} Biochemistry \& Biophysics Program, Amherst College, Amherst, MA
    }
    {
    This study attempts to verify a previously proposed hypothesis for histone-protamine replacement. When this replacement pathway goes wrong, it can lead to changes in the morphology, which is the second leading cause of male infertility. However, the physical dynamics of this pathway are currently unclear, limiting our ability to develop more pinpointed and effective treatments for male infertility. We hope that gaining a clearer understanding of this replacement pathway will enable further advancements in the field.
    }
    {
    To investigate the structural dynamics of this process, mononucleosomes were reconstituted using a single-tube salt-dilution protocol, a derivative of traditional salt dialysis, and incubated across a protamine titration (0.1{\textendash}1.0 \ensuremath{\mu}M). Atomic Force Microscopy (AFM) imaging was then used to create topographical images of our samples.
    }
    {
    Atomic Force Microscopy (AFM) imaging revealed that protamine promoted DNA looping and stabilized nucleosome architecture. No distinct $\sim$195 nm intermediate complex was observed. The absence of this intermediate validates previous lab findings supporting a rapid direct competition model over a stable, multi-step replacement pathway.
    }
    {
    While structural transition states were not observed during protamine titration, the results provide some validation for our lab{\textquoteright}s working framework of chromatin remodeling. Notably, the absence of a $\sim$195 nm intermediate complex supports a direct competition model over a screening model, confirming findings from previous work done in the lab by Emily Ma {\textquoteleft}20. In addition, protamine addition was confirmed to consistently promote DNA looping. From the analysis of the data, it seems that the addition to protamine may also help nucleosome formation. Currently, we do not know the implications of this. Taken together, these observations successfully replicate and reinforce our existing lab models, indicating that protamine facilitates chromatin reorganization by inducing structural DNA looping while competing with histone octamers for DNA binding. Moving forward, next steps will focus on troubleshooting reaction conditions. Specifically, we will need to optimize protamine-nucleosome incubation times and explore a broader range of protamine concentrations to better capture transient reaction dynamics.
    }

    \newpage
    